% ==============================================================================
% CHAPITRE 6 : TESTS, VALIDATION & BENCHMARKS DE PERFORMANCE
% ==============================================================================

\chapter{Tests, Validation et Benchmarks de Performance}
\label{ch:experiments}

\section{Environnement expérimental et jeu de données IBM AMLSim}

L'évaluation a été menée sur une machine hôte 8 cœurs (3.2 GHz), 16 Go de RAM sous Windows 11 / WSL2 Debian. Les conteneurs Docker allouent 4 cœurs CPU et 4 Go de RAM à Apache Spark 3.5.1, et 2 Go à Apache Kafka 3.5.1. Le jeu de données volumétrique de référence \textbf{IBM AMLSim} \cite{weber2018scalable} (\texttt{20K\_cycle200}) comprend $\mathbf{20\,000}$ comptes actifs, $\mathbf{117\,805}$ transactions horodatées et $\mathbf{200}$ cycles de blanchiment injectés représentant la vérité terrain. Aucun attribut de fraude n'est injecté dans le flux (\textit{Zero Label Leakage}).

\section{Matrice de conformité des 22 tests d'acceptation (AT-01 à AT-22)}

Une suite complète de 22 tests d'acceptation a été automatisée avec \texttt{pytest}. Tous les tests ont été validés avec un taux de succès de $100\,\%$, comme résumé dans le Tableau~\ref{tab:acceptance_tests}.

\begin{table}[htbp]
\centering
\caption{Validation de la matrice des 22 critères d'acceptation fonctionnelle}
\label{tab:acceptance_tests}
\scriptsize
\setlength{\tabcolsep}{2.5pt}
\renewcommand{\arraystretch}{0.90}
\begin{tabularx}{\textwidth}{c X c | c X c}
\toprule
\textbf{Code} & \textbf{Spécification Validée} & \textbf{Résultat} & \textbf{Code} & \textbf{Spécification Validée} & \textbf{Résultat} \\
\midrule
\textbf{AT-01} & Triangle simple intra-fenêtre & \textbf{\color{emeraldSuccess}PASS} & \textbf{AT-12} & Virement conflictuel (ID identique) & \textbf{\color{emeraldSuccess}PASS} \\
\textbf{AT-02} & Invariance par permutation rotatoire & \textbf{\color{emeraldSuccess}PASS} & \textbf{AT-13} & JSON malformé / Décimales & \textbf{\color{emeraldSuccess}PASS} \\
\textbf{AT-03} & Arrivée multi-lots d'arêtes & \textbf{\color{emeraldSuccess}PASS} & \textbf{AT-14} & Horodatage futur abusif & \textbf{\color{emeraldSuccess}PASS} \\
\textbf{AT-04} & Chaîne ouverte ($A \to B \to C$) sans cycle & \textbf{\color{emeraldSuccess}PASS} & \textbf{AT-15} & Devises mixtes dans le cycle & \textbf{\color{emeraldSuccess}PASS} \\
\textbf{AT-05} & Auto-virement ($A \to A$) & \textbf{\color{emeraldSuccess}PASS} & \textbf{AT-16} & Cycle d'ordre 4 ($A \to \dots \to A$) & \textbf{\color{emeraldSuccess}PASS} \\
\textbf{AT-06} & Événements distants ($> 10\text{ min}$) & \textbf{\color{emeraldSuccess}PASS} & \textbf{AT-17} & Carrefour marchand légitime & \textbf{\color{emeraldSuccess}PASS} \\
\textbf{AT-07} & Bornes d'intervalles $[start, end)$ & \textbf{\color{emeraldSuccess}PASS} & \textbf{AT-18} & Reconstruction hors ligne Delta & \textbf{\color{emeraldSuccess}PASS} \\
\textbf{AT-08} & Déduplication fenêtres imbriquées & \textbf{\color{emeraldSuccess}PASS} & \textbf{AT-19} & Reprise après crash et arrêt & \textbf{\color{emeraldSuccess}PASS} \\
\textbf{AT-09} & Arrivée désordonnée intra-fenêtre & \textbf{\color{emeraldSuccess}PASS} & \textbf{AT-20} & Tolérance pannes publication & \textbf{\color{emeraldSuccess}PASS} \\
\textbf{AT-10} & Antériorité au Watermark & \textbf{\color{emeraldSuccess}PASS} & \textbf{AT-21} & Non-avancement au repos & \textbf{\color{emeraldSuccess}PASS} \\
\textbf{AT-11} & Transaction dupliquée & \textbf{\color{emeraldSuccess}PASS} & \textbf{AT-22} & Profils de charge L1, L2, L3 & \textbf{\color{emeraldSuccess}PASS} \\
\bottomrule
\end{tabularx}
\end{table}

\section{Résultats des benchmarks de performance sous charge (AT-22)}

Le banc d'essai standardisé (\texttt{scripts/run\_benchmarks.py}) a évalué le pipeline sous trois paliers de débit : L1 Baseline (10 TPS), L2 Standard (100 TPS) et L3 Stress (500 TPS). Les résultats sont présentés dans le Tableau~\ref{tab:benchmarks}.

\begin{table}[htbp]
\centering
\caption{Résultats des campagnes de benchmarks sous charge d'ingestion (AT-22)}
\label{tab:benchmarks}
\vspace{4pt}
\scriptsize
\setlength{\tabcolsep}{4pt}
\renewcommand{\arraystretch}{1.08}
\begin{tabularx}{\textwidth}{l *{6}{>{\centering\arraybackslash}X}}
\toprule
\textbf{Palier} & \makecell{\textbf{Débit}\\\textbf{Cible}} & \makecell{\textbf{Volume}\\\textbf{Injecté}} & \makecell{\textbf{Latence}\\\textbf{P50}} & \makecell{\textbf{Latence}\\\textbf{P95}} & \textbf{Rappel} & \makecell{\textbf{Mémoire}\\\textbf{RAM}} \\
\midrule
\textbf{Niveau 1 (L1)} & $10\text{ TPS}$ & $106\text{ txns}$ & $0.43\text{ ms}$ & $0.43\text{ ms}$ & $\mathbf{100\,\%}$ & $170.8\text{ Mo}$ \\
\textbf{Niveau 2 (L2)} & $100\text{ TPS}$ & $515\text{ txns}$ & $1.47\text{ ms}$ & $1.47\text{ ms}$ & $\mathbf{100\,\%}$ & $199.1\text{ Mo}$ \\
\textbf{Niveau 3 (L3)} & $500\text{ TPS}$ & $1\,515\text{ txns}$ & $18.95\text{ ms}$ & $\mathbf{18.95\text{ ms}}$ & $\mathbf{100\,\%}$ & $227.4\text{ Mo}$ \\
\bottomrule
\end{tabularx}
\end{table}

\screenshotpaired[2.4cm]%
{figures/fig_benchmark_throughput_latency.png}{Latence P95 vs. Débit (L1, L2, L3)}%
{figures/fig_pagerank_distribution.png}{Distribution PageRank (Isolation Marchande AT-17)}%
{Évaluation des performances de latence et validation spectrale de non-régression}%
{fig:benchmarks_and_pagerank}

\subsection{Analyse des résultats obtenus}

Trois constats ressortent des mesures : 1) \textbf{Latence sub-seconde ($P95 = 18.95\text{ ms}$)} : même à $500\text{ TPS}$, le traitement prend moins de $20\text{ ms} \ll 30\,000\text{ ms}$ (pas de glissement), sans retard de file ; 2) \textbf{Rappel parfait ($100\,\%$)} : tous les cycles injectés sont isolés sans faux négatifs ; 3) \textbf{Empreinte mémoire contenue ($< 230\text{ Mo}$)} grâce à l'éviction continue opérée par le watermark.

\newpage
